% !TEX program = lualatex
\documentclass[11pt]{ltjsarticle}
\usepackage[margin=1in]{geometry}

% Core packages
\usepackage{amsmath,amssymb}
\usepackage{tikz-cd}
\usepackage{multicol}
\usepackage{url}

% Paragraphs
\setlength{\parindent}{0pt}
\setlength{\parskip}{1\baselineskip}

\title{トポス理論における古典論理の埋め込みと真理保存性の依存性：\\
メビウス的位相構造による連続体仮説の解釈}
\author{千葉将臣}
\date{2026-07-08}

\begin{document}
\maketitle

\begin{abstract}
本稿では、M. Tierney による層理論を用いた連続体仮説（continuum hypothesis, CH）の独立性の解釈を、古典論理の直観主義論理への埋め込みという観点から再検討する。とくに、前層トポスの直観主義的内部論理において、二重否定位相を通じて古典論理が実現される過程に注目する。この過程は、古典論理が直観主義論理へ埋め込めることを示す一方で、埋め込まれた古典論理における真理保存性が、背後の直観主義的基盤、サイトの選択、および層化の方式に依存して現れることを示唆する。本稿はこの依存性を、メビウス帯における局所的な表裏の区別と大域的な同一性の類比によって説明し、CH の独立性を、真偽の単なる未決定性ではなく、真理保存性の相対化と不可視化の構造として解釈する。
\end{abstract}

\section{導入：層の公理的理論による構成の内在化}

古典的な強制法では、集合論的宇宙の外部からジェネリック・フィルターを選択し、拡大モデルを構成するというメタレベルの操作が重要な役割を果たす。これに対して、Lawvere と Tierney によって展開された層およびトポスの方法は、この外部的操作を、対象レベルの構造として内在化する道を開いた。Tierney の「層理論と連続体仮説」において中心となるのは、集合論的独立性の議論を、トポスの内部論理および位相の選択によって再構成するという視点である\cite{tierney1972}。

本稿の関心は、CH の独立性証明そのものを再証明することではない。むしろ、Tierney 的構成において、古典論理が直観主義的なトポスの内部にどのように実現されるのか、そしてその実現後の古典論理における真理保存性が、どのような意味で背後の構成方式に依存するのかを問うことにある。この問題意識は、形式体系、証明、意味論、および真理の関係を明示的に扱う Kleene のメタ数学的精神を引き継ぐものである\cite{kleene1952}。

現代の集合論的実践では、しばしば十分強いメタ理論を固定し、その内部で対象理論やモデルを扱う。その場合、推論や翻訳が真理保存的であるかという問いは、背景的条件として処理されることが多い。しかし、古典論理を直観主義的基盤の中に埋め込む場面では、この問いは再び可視化される。すなわち、問題は「古典論理を直観主義論理へ埋め込めるか」ではない。むしろ、埋め込まれた古典論理において成立する真理保存性が、どの基盤から、どの位相を通じて、どのように得られたものなのかという点にある。

\section{二重否定位相と古典論理の実現}

前層トポスの内部論理は一般に直観主義論理である。そこでは命題の真理値は単なる真と偽の二値に限られず、部分対象分類子によって表される豊かな構造をもつ。このため、ある命題 $A$ とその二重否定 $\neg\neg A$ は、古典論理においては同値であっても、直観主義論理においては一般には区別される。

古典論理をこの直観主義的環境の中に実現する標準的な方法の一つが、二重否定位相である。二重否定位相は、直観主義的な真理値構造に対して、$A$ と $\neg\neg A$ の差異を層化によって捨象し、排中律が成立するブール的な環境を取り出す。したがって、古典論理は直観主義論理と無関係に外部から付加されるのではなく、直観主義的トポスの内部で、特定の位相を通じて実現される。

この点を本稿では、メビウス帯の初等的比喩によって捉える。メビウス帯では、局所的には表と裏を区別できるように見える。しかし、大域的にはその二つは一続きの面として接続される。同様に、前層段階では $A$ と $\neg\neg A$、あるいは未確定な真理状態と古典的真理状態は区別されるが、二重否定位相による層化を経ると、それらの差異は古典的な真理構造の中で不可視化される。

ただし、この比喩は技術的同一視そのものではなく、解釈上の補助線である。メビウス的構造という語は、真と偽、表と裏、直観主義的未確定性と古典的決定性が、単純に対立する二領域としてではなく、ある変換を通じて接続されるという意味で用いる。この発想は、真理保存性を位相幾何学的モデルによって説明しようとする千葉の議論から示唆を受けている\cite{chiba2026}。

\section{CH の独立性と埋め込まれた古典論理の盲点}

CH の独立性は、古典的集合論において、CH を満たすモデルと $\neg\mathrm{CH}$ を満たすモデルの双方が相対的一貫性の範囲で構成されることによって示される。Tierney 的な層・トポスの観点では、この独立性は、サイト、位相、層化の選択に応じて、内部論理における命題の成立条件が変化する現象として再解釈される。

ここで重要なのは、CH の独立性が単に「真でも偽でもよい」という任意性を意味するのではないという点である。むしろ、直観主義的基盤の段階では、CH を真とする構成と $\neg\mathrm{CH}$ を真とする構成が、それぞれ異なる仕方で与えられうる。そして二重否定位相を通じて古典論理が実現されると、得られた古典的内部論理の側からは、その古典的宇宙がどの直観主義的基盤や位相選択から生じたのかを一意に回収することができない。

この不可回収性を、本稿では「埋め込まれた古典論理の盲点」と呼ぶ。外部のメタ理論からは、異なる直観主義的基盤、異なるサイト、異なる層化の方式が見えている。しかし、埋め込まれた古典論理の内部では、排中律と二値性が成立するため、その背後にあった直観主義的な差異は見えなくなる。メビウス帯の局所的観察者が、自分のいる面の大域的な反転構造をただちには把握できないのと同様に、古典論理の内部観察者は、自らの真理構造を生成した基盤の差異を直接には対象化できない。

したがって、CH の独立性は、古典論理の内部で単に証明不能な命題が存在するという事実にとどまらない。それは、古典論理が直観主義的基盤上に実現されたとき、その内部からは見えない構成依存性を伴うという構造的事実として読める。この意味で、CH と $\neg\mathrm{CH}$ は、互いに断絶した二つの世界に属するのではなく、層化と埋め込みを通じて、同一の意味論的表面上にメビウス的に現れる。

\section{真理保存性の同型的拡張}

通常、真理保存性とは、ある推論、翻訳、または埋め込みが、真である命題を真である命題へ写すことを意味する。しかし、トポス的文脈では、真理値は単なる二値ではなく、部分対象分類子や位相によって構造化された対象として現れる。このため、真理保存性を命題ごとの真偽の一致だけに限定すると、層化や内部言語の変換において実際に保存されているものを十分に捉えられない。

そこで本稿では、真理保存性を同型的に拡張して理解する。すなわち、真理保存性とは、個々の命題の真偽をそのまま保つことだけではなく、真理値が配置される意味論的構造の対応関係を保つことでもある。この拡張された意味では、ある体系から別の体系への翻訳が、真理構造を構造保存的に対応させるとき、その翻訳は同型的な真理保存性をもつと言える。

しかし、この同型的保存性は、体系内部からは必ずしも可視的ではない。外部から見れば、直観主義的基盤と古典的層化との間には構造的対応が存在する。しかし、埋め込まれた古典論理の内部では、その対応全体を対象化するための視点が失われている。ここに、真理保存性の同型的拡張がもたらす盲点がある。

この盲点は、真理保存性が失敗しているというより、真理保存性の成立条件が内部化された結果、その条件そのものが内部からは見えなくなるという事態である。したがって、埋め込まれた古典論理における真理保存性は、体系単独で自足する性質ではなく、背後の直観主義的基盤と層化の方式に依存して発生する性質である。

\section{結論：埋め込まれた古典論理における真理保存性の依存性}

本稿の主張は、古典論理が直観主義論理へ埋め込めないということではない。むしろ、二重否定位相を通じて、古典論理は直観主義的トポスの内部に実現されうる。本稿が問題にしたのは、その実現後の古典論理において成立する真理保存性の身分である。

二重否定位相による層化は、直観主義的な中間状態を捨象し、古典的なブール的環境を構成する。しかし、そのことは、埋め込まれた古典論理の真理概念が、背後の直観主義的基盤から独立して自足することを意味しない。むしろ、古典論理の内部で見える真理は、どのサイトを選び、どの位相を適用し、どの層化を経てその古典的環境が得られたのかに相対化されている。

したがって、古典論理は直観主義論理に埋め込めるが、埋め込まれた古典論理には真理保存性への依存が発生する。この依存性は、古典論理の内部からは盲点として現れる。CH の独立性は、この盲点をもっとも鮮明に示す例である。すなわち、CH と $\neg\mathrm{CH}$ は、古典的内部論理においては決定不能性として現れるが、外部のトポス的・層的観点からは、異なる基盤と位相選択に相対化された真理構造として理解される。

この意味で、真理保存性は、古典的体系単独で完結する絶対的性質ではなく、埋め込み、層化、位相選択を通じて生成される構造的性質である。本稿が「メビウス的」と呼ぶのは、この構造である。すなわち、局所的には真と偽、表と裏、古典と直観主義が区別されるにもかかわらず、大域的にはそれらが一つの意味論的構造の中で反転しつつ接続されるという構造である。

\begin{thebibliography}{9}

\bibitem{kleene1952}
Kleene, S. C. (1952).
\newblock {\itshape Introduction to Metamathematics}.
\newblock North-Holland.

\bibitem{tierney1972}
Tierney, M. (1972).
\newblock Sheaf theory and the continuum hypothesis.
\newblock In D. Bucur, A. Deleanu, and P. J. Hilton (eds.), {\itshape Toposes, Algebraic Geometry and Logic}, Lecture Notes in Mathematics, Vol. 274, Springer-Verlag, pp. 13--42.

\bibitem{chiba2026}
千葉将臣 (2026).
\newblock 真理保存性の不完全性定理の位相幾何学的モデル：メビウス多様体における自己言及構造の解消.
\newblock \url{https://purl.org/aletheics/papers/sunyata-math/titotp-mebius.pdf}.

\end{thebibliography}

\end{document}
